\section[Aplicación a los morfismos finitos y cuasifinitos]{Aplicación a los morfismos finitos y cuasifinitos\protect\footnotemark}
\label{VIII.6}
\footnotetext{\Cf EGA IV 18.12 para generalizaciones a preesquemas no necesariamente localmente noetherianos}

Vamos a demostrar los dos teoremas siguientes:

\begin{theoreme}
\label{VIII.6.1}
Sea~$f\colon X\to Y$ un morfismo \emph{propio con fibras finitas}, con~$Y$
localmente noetheriano. Entonces~$f$ es finito.
\end{theoreme}

\begin{theoreme}
\label{VIII.6.2}
Sea~$f\colon X\to Y$ un morfismo \emph{cuasifinito} y \emph{separado},
con~$Y$ localmente noetheriano. Entonces~$f$ es cuasiafín y, a fortiori,
cuasiproyectivo.
\end{theoreme}

\begin{remarques}
\label{VIII.6.3}
El
\marginpar{217}
teorema~\Ref{VIII.6.1} es bien conocido y se debe a Chevalley en el caso de
las variedades algebraicas; también se encontrará una demostración sencilla
en [EGA III 4], que utiliza el «teorema de las funciones holomorfas». La
demostración dada aquí no utiliza este último teorema, sino la teoría del
descenso; la damos como obsequio al lector, pues se obtiene «gratis» al mismo
tiempo que la de~\Ref{VIII.6.2}. Recordemos también ([EGA III 4] o
\cite{VIII.1}) que la forma global del \emph{Main Theorem} de Zariski,
deducida del «teorema de las funciones holomorfas», afirma que, si
$f\colon X\to Y$ es cuasifinito y \emph{cuasiproyectivo}, siendo~$Y$
noetheriano, entonces~$X$ es $Y$-isomorfo a un subpreesquema abierto de un
$Y$-preesquema \emph{finito}~$Z$. La conjunción del \emph{Main Theorem} y
de~\Ref{VIII.6.2} se enuncia, por tanto, así:
\end{remarques}

\begin{corollaire}
\label{VIII.6.4}
Sea~$f\colon X\to Y$ un morfismo cuasifinito y separado, con~$Y$
noetheriano. Entonces~$X$ es $Y$-isomorfo a un subpreesquema abierto de un
$Y$-preesquema finito~$Z$.
\end{corollaire}

Otra consecuencia interesante de~\Ref{VIII.6.2} para la teoría del descenso
se dará junto con~\Ref{VIII.7.9}.

\subsubsection*{Demostración de~\Ref{VIII.6.1} y~\Ref{VIII.6.2}}
Admitiremos el hecho siguiente, cuya demostración es
fácil\footnote{\Cf EGA IV 8}:

\begin{lemme}
\label{VIII.6.5}
Sea~$X$ un preesquema de tipo finito sobre~$Y$ localmente noetheriano, y
sea~$y\in Y$. Para que exista un entorno abierto~$U$ de~$y$ tal que~$X|U$
sea finito (\resp cuasiafín, \resp\ ...) sobre~$U$, es necesario y suficiente
que~$X\times_{Y}\Spec(\mathcal{O}_{y})$ sea finito (\resp cuasiafín,
\resp\ ...) sobre~$\Spec(\mathcal{O}_{y})$.
\end{lemme}

Como, por otra parte, la propiedad de~$f\colon X\to Y$ de ser finito,
\resp cuasiafín, es local sobre~$Y$, para demostrar~\Ref{VIII.6.1}
y~\Ref{VIII.6.2} se reduce uno al caso en que~$Y$ es el espectro de un anillo
local y, a fortiori, tiene dimensión finita. (N.B.: se llama
\emph{dimensión de un preesquema}~$Y$ al supremo de las dimensiones de Krull
de sus anillos locales). Procedemos por inducción sobre
$$
n=\dim(Y)\quoi,
$$
siendo la afirmación
\marginpar{218}
trivial para~$n<0$. Podemos suponer, por tanto, $n\geq 0$ y demostrada la
afirmación para las dimensiones~$n'<n$. Puede suponerse de nuevo que~$Y$ es
el espectro de un anillo local noetheriano~$A$, de dimensión~$n$. Obsérvese
que las hipótesis hechas en~\Ref{VIII.6.1} y~\Ref{VIII.6.2} son estables por
cambio de base (ya se ha hecho uso de ello en la reducción inicial); seguirán
siendo ciertas tras el cambio de base~$\Spec(\widehat{A})\to \Spec(A)$. Como
este último es fielmente plano y cuasicompacto, los
enunciados~\Ref{VIII.5.7} y~\Ref{VIII.5.9} reducen la cuestión al caso en
que~$A$ es, además, completo. Utilizando entonces el hecho de que todo anillo
local noetheriano~$B$ sobre~$A$ que es cuasifinito sobre~$A$ es finito
sobre~$A$, y el hecho de que~$X$ es separado sobre~$Y$ y la fibra sobre~$y$ está formada por
puntos aislados, se obtiene una descomposición
$$
X=X'\amalg X''
$$
en la que~$X'$ es \emph{finito} sobre~$Y$ y la fibra de~$X''$ en~$y$ es
vacía. Si~$X$ es propio sobre~$Y$, lo mismo sucede con~$X''$; su imagen en
$Y$ es, por tanto, cerrada y, como no contiene a~$y$, es vacía. Así,
$X''=\emptyset$ y, por consiguiente, $X=X'$, lo que demuestra que~$X$ es
finito sobre~$Y$ y prueba~\Ref{VIII.6.1} (N.B.: aquí es inútil la hipótesis
de inducción). Si~$X$ es cuasifinito sobre~$Y$, $X''$ también lo es; ahora
bien, $X''$ se encuentra de hecho sobre
\ifthenelse{\boolean{orig}}
{el abierto~$Y-(y)$}
{el abierto~$Y-\{ y\}$}
de~$Y$, \emph{que tiene dimensión} $<n$. En virtud de la hipótesis de
inducción, $X''$ es cuasiafín
\ifthenelse{\boolean{orig}}
{sobre~$Y-(y)$,}
{sobre~$Y-\{ y\}$,}
y, por tanto, también sobre~$Y$; evidentemente, lo mismo sucede con~$X'$ y,
por consiguiente, con su suma~$X$, lo que prueba~\Ref{VIII.6.2}.

\begin{remarque}
\label{VIII.6.6}
Los teoremas~\Ref{VIII.6.1} y~\Ref{VIII.6.2} siguen siendo válidos si ya no
se supone~$Y$ localmente noetheriano, con tal de precisar que se supone~$f$
de presentación finita
\ifthenelse{\boolean{orig}}
{(\cf \Ref{VIII.3.6}).}
{(\cf \Ref{VIII.3.6}).}
En efecto, todavía puede suponerse~$Y$ afín; se comprueba entonces sin
dificultad que la situación~$f\colon X\to Y$
\ifthenelse{\boolean{orig}}
{se deduce,}
{se deduce,}
mediante un cambio de base~$Y\to Y_{0}$, de una situación
$f_{0}\colon X_{0}\to Y_{0}$ que satisface las mismas hipótesis que~$f$,
con~$Y_{0}$ \emph{noetheriano}. Por tanto, según el
resultado~\Ref{VIII.6.1} \resp \Ref{VIII.6.2}, $f_{0}$ es finito
\resp cuasiafín y, en consecuencia, lo mismo sucede con~$f$. Este género de
razonamiento resulta útil a menudo para librarse de hipótesis noetherianas
\marginpar{219}
(que siempre acaban siendo molestas en las aplicaciones).
\end{remarque}
